\section{Automatic norm certificates and sharp finite zero counts}
\label{sec:norm-certificates}

The earlier embedding packet proves a sharp zero bound for
$y^2=f(x)$ when $\deg f$ is odd, using
$A^2-fB^2$. It also gives a general finite-algebra spanning-certificate
interface. The remaining question is how to \emph{construct} useful
certificates for wider classes without an unspecified elimination
algorithm \cite{repo-embedding,repo-siegel}. The present section supplies such a constructor for every
monic equation in one coordinate, and an automatic nonvanishing proof
for the coprime superelliptic family. For this family it replaces the
divisor-based zero-count obligation in the Siegel packet by explicit
polynomial and matrix arithmetic.

All fields used below have exact, decidable equality and primitive
recursive arithmetic on finite codes, with their field laws available
in the object theory. Rational numbers are the principal example. A
fixed explicitly presented number field works as well once those
arithmetic contracts have been proved. Arbitrary Cauchy names for real
or complex coefficients do not provide this interface. Lists of points
are always finite input data; no set of all points is formed.

\subsection{The determinant constructor for a monic plane equation}

Let $K$ be such a coded field and let
\[
 F(X,Y)=Y^m+\sum_{i<m}f_i(X)Y^i,\qquad m\geq1.
\]
Every polynomial in $Y$ has a remainder of degree less than $m$ after
division by $F$, computed by monic division. Write
$G=\sum_{i<m}A_i(X)Y^i$ for the remainder of the function whose zeros
are to be counted. For each $j<m$, reduce $GY^j$ modulo $F$ and put
its coefficient vector in column $j$ of an $m$ by $m$ matrix $M_G(X)$.
Define
\[
                         N_G(X)=\det M_G(X).
\]
The matrix and determinant are finite polynomial computations. In
particular this definition does not presume an algebraic closure,
factorization, a divisor, or a terminating search for a certificate.

\begin{theorem}[Finite norm zero bound]\label{thm:general-norm}
If $N_G\ne0$, every finite list of pairwise distinct points
$(x,y)\in K^2$ satisfying $F(x,y)=G(x,y)=0$ has length at most
$\deg N_G$. Nonvanishing of the computed polynomial is a finite
coefficient test. The statement includes singular curves and
ramified fibers.
\end{theorem}

\begin{proof}
First recall the elementary finite root bound: a nonzero polynomial of
degree $r$ over a field has at most $r$ distinct roots in a supplied
list. Synthetic division removes $Y-b$ at a listed root $b$. At a
different listed root $c$, its nonzero value $c-b$ implies that the
quotient still vanishes. A finite induction removes the product of
the distinct linear factors, whose degree cannot exceed $r$. This
is a proof by polynomial operations on the input list.

Group the supplied common zeros by their $x$-coordinates. Fix one
group with common abscissa $a$ and distinct ordinates
$b_1,\ldots,b_t$. The preceding root bound, applied to the monic
polynomial $F(a,Y)$, gives $t\leq m$. Set
\[
               e(b)=(1,b,\ldots,b^{m-1}).
\]
The definition of $M_G$ and evaluation of the monic-division
identities give
\[
 e(b_\ell)M_G(a)=G(a,b_\ell)e(b_\ell)=0
                    \qquad(1\leq\ell\leq t).
\]
Extend these $t$ rows to a square matrix $U$ by adjoining the
standard rows with indices $t,t+1,\ldots,m-1$, using zero-based
column indices. Its upper-left $t$ by $t$ block is the Vandermonde
matrix on the distinct $b_\ell$, and its lower-left block is zero
with lower-right block the identity. Thus
\[
             \det U=\prod_{1\leq i<j\leq t}(b_j-b_i)\ne0.
\]
The Vandermonde identity itself has a finite inductive proof.
Replace the last ordinate by an indeterminate $T$. The determinant
has degree at most $t-1$ in $T$ and vanishes at the preceding
$t-1$ distinct ordinates because two rows then agree. The root
removal argument already proved above divides out all
$T-b_i$. Its leading coefficient is the preceding $(t-1)$ by
$(t-1)$ Vandermonde determinant, by expansion in the last row.
Induction, starting with the determinant one in size one, proves
the formula after substitution $T=b_t$.

Every entry in the first $t$ rows of $UM_G(X)$ vanishes at $X=a$.
Synthetic division therefore removes a factor $X-a$ from each such
entry. Multilinearity in the rows shows
\[
                        (X-a)^t\mid\det(UM_G(X)).
\]
Since $\det U$ is a nonzero scalar and
$\det(UM_G)=\det U\,N_G$, it follows that
$(X-a)^t\mid N_G$. This step counts the \emph{number of distinct
points} in the fiber, not merely the number of abscissae. In
particular a ramified point with $b=0$ causes no division by $b$;
all divisions in $U^{-1}$, if one chooses to compute it, are by the
witnessed nonzero differences $b_j-b_i$ and by a nonzero determinant.
The proof does not need to compute that inverse.

For different abscissae, multiply these divisibilities. Here is an
explicit cancellation argument that avoids a factorization oracle.
Suppose $N_G=CQ$ after removing the factors belonging to previous
abscissae. At a new abscissa $a$ we have $C(a)\ne0$. From
$(X-a)^t\mid CQ$ and $t>0$, evaluation gives $Q(a)=0$. Write
$Q=(X-a)Q_1$ by synthetic division and cancel one $X-a$ from the
two polynomial expressions. Repeat exactly $t$ times. This removes
the new factor without disturbing the previous ones.

Consequently, if the groups have sizes $t_a$, then
\[
               \prod_a(X-a)^{t_a}\mid N_G.
\]
The product is monic of degree $\sum_a t_a$, which equals the
original list length. The assumption $N_G\ne0$ now proves the
required bound.
\end{proof}

There is also a fully explicit coarse degree estimate. Suppose all
nonzero $A_i$ have degree at most $a$ and all nonzero $f_i$ have
degree at most $c$, where $a,c\geq0$. The multiplication matrix for
$Y$ has entries of degree at most $c$: its nontrivial last column
is $(-f_0,\ldots,-f_{m-1})^T$. Therefore
$M_G=\sum_{i<m}A_iT^i$ has entries of degree at most
$a+(m-1)c$. Expanding its determinant gives
\[
              \deg N_G\leq m\bigl(a+(m-1)c\bigr)
                       \quad\text{if }N_G\ne0.
\]
The nonvanishing check is essential for a general monic equation.
For example $F=Y^2-X^2$ and $G=Y-X$ give $N_G=0$ and infinitely
many common rational zeros. The next theorem supplies a structural
nonvanishing proof for an important uniform class.

\subsection{Coprime superelliptic equations: no nonvanishing oracle}

Fix $m\geq2$ and a nonzero $f\in K[X]$ with positive degree $d$.
Assume
\[
                             \gcd(m,d)=1.
\]
For $G=\sum_{i<m}A_i(X)Y^i\ne0$, define its weighted degree by
\[
 W(G)=\max_{\substack{i<m\\A_i\ne0}}
                   \bigl(m\deg A_i+d i\bigr).
\]
The maximum has a unique maximizing index $k$. Indeed equality of
two displayed weights implies $m\mid d(i-j)$, hence $m\mid i-j$,
and $0\leq i,j<m$ then forces $i=j$. Coprime cancellation here can
be derived from the integer Bezout identity produced by bounded
Euclidean arithmetic.

\begin{theorem}[Exact degree of the superelliptic norm]
\label{thm:norm-degree}
For $F=Y^m-f(X)$, the determinant constructor above gives
\[
                       N_G\ne0,\qquad \deg N_G=W(G).
\]
If $k$ is the unique maximizing index, its leading coefficient is
\[
 \operatorname{sgn}(\sigma_k)
       \operatorname{lc}(A_k)^m\operatorname{lc}(f)^k,
 \qquad \sigma_k(j)=(j+k)\bmod m.
\]
\end{theorem}

\begin{proof}
Use row and column indices $r,j\in\{0,\ldots,m-1\}$. There is
exactly one $i$ in this range with $r\equiv i+j\pmod m$.
Writing $i+j=r+mq$, we have $q\in\{0,1\}$, and reduction of
$Y^{i+j}$ gives the particularly simple matrix entry
\[
                       (M_G)_{rj}=A_i f^q.
\]
Whenever this entry is nonzero,
\[
       m\deg (M_G)_{rj}+d r-d j
            =m\deg A_i+d i\leq W(G).
\]
Consider a nonzero determinant term associated to a permutation
$\sigma$. Sum this equality over its $m$ columns. The row and
column correction terms cancel because $\sigma$ permutes the
indices, so
\[
 m\deg\left(\prod_{j<m}(M_G)_{\sigma(j),j}\right)
       =\sum_{j<m}\bigl(m\deg A_{i_j}+d i_j\bigr)
       \leq mW(G).
\]
Equality is possible only if every $i_j=k$, by uniqueness of the
maximizing index. Exactly one permutation does this, namely
$\sigma_k(j)=(j+k)\bmod m$. Among these entries precisely $k$
wrap around the end of the index range. Their product is therefore
$A_k^m f^k$, of degree $m\deg A_k+dk=W(G)$.

Every other nonzero determinant term has strictly smaller degree.
Thus the top coefficient of the unique term cannot cancel. Its
value is the nonzero coefficient displayed in the statement, which
proves both assertions. This is a direct determinant computation;
no irreducibility theorem, valuation extension, or algebraic closure
is used.
\end{proof}

\begin{corollary}[Sharp weighted zero count]\label{cor:norm-zeros}
Every finite list of pairwise distinct points of $y^m=f(x)$ on
which $G(x,y)=0$ has length at most $W(G)$.
\end{corollary}
\begin{proof}
Apply Theorems~\ref{thm:general-norm} and~\ref{thm:norm-degree}.
\end{proof}

The degree hypothesis genuinely matters to the automatic proof.
When $m=d=2$, $f=X^2$ and $G=Y-X$, both summands have weight two;
their norm cancels identically. The general determinant theorem
remains valid whenever its computed norm is nonzero, even when
$\gcd(m,d)>1$.

The weighted bound is attained over $\Q$. Choose distinct rational
numbers $a_1,\ldots,a_d$ and put
\[
           f(X)=1+\prod_{j=1}^d(X-a_j),\qquad G=Y-1.
\]
For $\gcd(m,d)=1$, the common zeros are exactly
$(a_1,1),\ldots,(a_d,1)$, and $W(G)=d$. For example $m=3,d=4$
and $(a_j)=(0,1,2,4)$ give
$f=X^4-7X^3+14X^2-8X+1$. This example illustrates sharpness of
the arithmetic zero bound; smoothness is neither required nor
claimed by the theorem.

\subsection{A uniform proper-subspace fiber theorem}

For any $L\in\N$, let
\[
 \mathcal B_L=\{X^jY^i:0\leq i<m,\ j\geq0,\ mj+di\leq L\},
 \qquad n_L=|\mathcal B_L|.
\]
The constant monomial belongs to this list, so $n_L\geq1$.
Enumerate its entries lexicographically, and define
$\Phi_L(x,y)\in K^{n_L}$ by evaluating them in that order.
The list is obtained by the loops $i<m$ and $j\leq L$ followed by
a numerical test; no symbolic basis search occurs.

\begin{corollary}[Subspace fibers for superelliptic curves]
\label{cor:norm-subspace}
For every proper subspace $V\subsetneq K^{n_L}$ supplied by a
finite matrix of spanning rows of rank less than $n_L$, every
finite list of distinct points on $y^m=f(x)$ whose evaluation
vectors lie in $V$ has length at most $L$.
\end{corollary}
\begin{proof}
Bounded Gaussian elimination produces a nonzero annihilating
linear functional: choose a free column, assign its unknown one,
set the other free unknowns to zero, and solve the pivot equations
backwards. Its evaluation on $\Phi_L$ is a nonzero canonical
polynomial $G=\sum_{i<m}A_i(X)Y^i$ with $W(G)\leq L$.
Every point under consideration is a zero of $G$.
Corollary~\ref{cor:norm-zeros} gives the result.
\end{proof}

If a separately proved finite-tuple theorem covers the evaluation
vectors of each supplied list of points in a decidable class $S$
by at most $M$ such proper subspaces, then every distinct list
from $S$ has length at most $ML$. Assign each point to the first
covering subspace by bounded exact membership search; each of the
at most $M$ groups has size at most $L$. This last finite counting
step is the existing Siegel interface. The new contribution is an
automatic, sharp proof of its geometric fiber contract for all
coprime exponents $m,d$, as well as a computable norm test for a
general monic plane equation.

This result does not construct the quantitative approximation
cover, a complete list of rational or integral points, or a global
height cutoff. It therefore does not prove general Siegel or
Faltings. Nor does it impose smoothness or genus hypotheses that
would be needed for those larger theorems. It resolves the stated
finite zero-count dependency on a substantially wider class than
the previous odd-degree hyperelliptic case.

\subsection{Arithmetic compilation and Lean obligations}

For rational coefficients, use canonical signed numerators and
positive denominators, coefficient lists for univariate polynomials,
and lists of such lists for $G$. Monic reduction runs for the
input $Y$-degree. A companion matrix, finite matrix powers, and a
determinant expansion construct the norm. Permutations can be
enumerated by scanning all length-$m$ lists of indices less than
$m$ and retaining those with no repeated entries: the candidate
count is $m^m$. Thus even the deliberately inefficient determinant
algorithm is one fixed primitive recursive program, uniformly in
$m$. Euclidean and Gaussian elimination loops are bounded by
their input lengths or matrix dimensions. Exact coefficient size
may grow, but integer and rational arithmetic on codes remains
primitive recursive.

The human proof translates by quantifier-free loop invariants.
Monic division tracks the coefficient identity
$H=QF+R$ and the degree bound on its current remainder. Matrix
construction tracks, for every bounded column index, the
coefficient identity for $GY^j$ modulo $F$. The determinant
degree argument uses finite products, bounded permutation lists,
and a Boolean test for the unique cyclic shift. Grouping points
tracks that the accumulated groups partition the input list;
root removal tracks $N_G=CQ$ and the accumulated degree. All
coefficient equalities, bounded membership assertions, and
nonzero tests are Boolean folds on these finite data. The field
laws and the arbitrary supplied list remain external parameters
of these inductions. No induction formula asserts a universal
claim about all points of the curve.

The file \path{lean/NormTargets.lean} gives executable
coefficient-list definitions, a general monic-plane companion matrix,
monic reduction, the superelliptic matrix, determinants, weighted
degrees, finite-point checks, and named propositions for the main
theorems. It also defines weighted monomial evaluation vectors, row-span
witnesses, and nonzero annihilator certificates for the subspace result.
These propositions are
\emph{unproved targets}; compiling their definitions proves none
of their mathematical content. A complete formalization should
proceed in the following dependency order:
\begin{enumerate}
\item Prove normalization, convolution, evaluation, degree of
  nonzero products, and synthetic division for the chosen codes.
\item Prove the monic-reduction identities and the evaluation-row
  equation for the multiplication matrix.
\item Prove determinant multiplicativity, row multilinearity,
  the Vandermonde identity, and finite root removal; assemble
  Theorem~\ref{thm:general-norm}.
\item Prove the unique-weight lemma by Bezout cancellation,
  the entry weight equality, and the unique cyclic-shift term;
  assemble Theorem~\ref{thm:norm-degree}.
\item Prove the finite Gaussian annihilator construction and
  use the existing finite-cover counting argument.
\item Reify the fixed arithmetic programs and each displayed
  bounded invariant in the repository's primitive-recursive
  syntax, and build the corresponding embedded
  \texttt{Prf PRAAx} derivations. Native Lean matrix theorems alone
  are not this final object-theory step.
\end{enumerate}

The reference program \path{code/superelliptic_norm.py}
constructs these matrices and determinants with exact rational
polynomial arithmetic. It checks leading degrees against the
formula, determinant evaluation identities, the hyperelliptic
specialization, a sharp cubic-cover example, and repeated factors
coming from two-point and ramified fibers. These are finite
implementation checks supporting the definitions, separate from
the all-input human proofs above.
